\cvsection{Профиль}
Системный аналитик с 5+ годами опыта на стыке бизнес-требований, интеграций,
данных и production-инфраструктуры. Специализируюсь на MDM-платформах,
микросервисных системах и API-интеграциях: превращаю бизнес-задачи в
понятные спецификации, контракты, модели данных и задачи для разработки.
Сильная сторона - умение говорить на одном языке с заказчиком, архитекторами,
разработкой, QA и эксплуатацией: от сбора требований и спецификаций до приемки
и готовности к промышленной эксплуатации.

\cvsection{Желаемая работа}
Интересны сложные высоконагруженные продукты и платформенные команды, где можно
глубоко работать с системным анализом, быть ближе к архитектурным решениям и
расти в сторону solution architecture. Готов к релокации.

\cvsection{Ключевые компетенции}
\begin{itemize}[leftmargin=1.3em, itemsep=0.2em]
  \item Аналитика: сбор и формализация требований, User Story, Use Case, бизнес-правила, сценарии валидации и обработки ошибок
  \item Интеграции: REST API, gRPC, Kafka, RabbitMQ, Swagger/OpenAPI, контракты API, форматы обмена JSON
  \item Данные: логическая модель данных, master data, сущности, атрибуты, связи, SQL, PostgreSQL, Redis, OpenSearch
  \item Процессы: декомпозиция требований, постановка задач разработчикам, согласование решений, демо заказчику, приемка
  \item Delivery: Jira, Confluence, Draw.io, Postman, Agile/Scrum, поддержка QA и анализ дефектов
  \item Инфраструктурный бэкграунд: Linux, Docker, Kubernetes/OpenShift, GitLab CI/CD, Jenkins, Ansible, Nginx, Prometheus, Grafana, ELK
\end{itemize}

\cvsection{Опыт работы}
\cvrole{Системный аналитик}{09.2023 -- 03.2026}{АО Консалт Плюс -- Фрегат}{Москва}
\textit{Фрегат: корпоративная MDM-платформа для централизованного управления мастер-данными и интеграции внутренних информационных систем.}
\begin{itemize}[leftmargin=1.3em, itemsep=0.2em]
  \item Собирал, анализировал и формализовал бизнес-требования совместно с заказчиком; переводил верхнеуровневые ожидания в проверяемые функциональные и нефункциональные требования.
  \item Декомпозировал требования в задачи для команды разработки, описывал бизнес-процессы, пользовательские сценарии и правила обработки данных.
  \item Проектировал интеграционные взаимодействия между внутренними сервисами, готовил спецификации для REST API, gRPC и Kafka.
  \item Описывал API-контракты, форматы обмена, структуры запросов и ответов, сценарии ошибок и правила валидации.
  \item Проектировал логическую модель данных MDM: сущности, атрибуты, связи, справочники и правила обработки мастер-данных.
  \item Готовил UML Sequence Diagram и схемы взаимодействия сервисов в Draw.io для согласования с заказчиком, архитекторами и разработкой.
  \item Консультировал разработчиков в процессе реализации, поддерживал QA при подготовке тестовых сценариев, анализировал дефекты и уточнял требования по итогам тестирования.
  \item Проводил демонстрации реализованной функциональности заказчику и участвовал в приемке перед промышленной эксплуатацией.
\end{itemize}

\cvrole{DevOps-инженер}{09.2019 -- 09.2022}{Сбер}{Москва}
\textit{Корпоративная система "Сенат" для поддержки управленческих решений руководителей банка. Команда проекта - более 30 человек.}
\begin{itemize}[leftmargin=1.3em, itemsep=0.2em]
  \item Развертывал и сопровождал высоконагруженные сервисы в банковской среде, администрировал Linux/Windows-серверы.
  \item Настраивал CI/CD-процессы и автоматизировал доставку приложений, работал с контейнеризацией и оркестрацией.
  \item Поддерживал мониторинг, отказоустойчивость и эксплуатационную готовность сервисов.
  \item Помогал командам разработки в инфраструктурных вопросах, что усилило системное понимание жизненного цикла требований от анализа до production.
\end{itemize}

\cvsection{Образование}
Донской государственный технический университет, Ростов-на-Дону\\
Бакалавр, Автоматизация, мехатроника и управление, 2021

\cvsection{Языки}
Русский - родной
